\documentclass[UTF8,zihao=-4]{ctexart}
\usepackage[a4paper,margin=2.2cm]{geometry}
\usepackage{fontspec}
\usepackage{tabularx}
\usepackage{array}
\usepackage{ragged2e}
\usepackage{fancyhdr}
\usepackage{hyperref}
\setCJKmainfont{Noto Serif SC}
\setCJKsansfont{Noto Sans SC}
\setmainfont{Noto Serif SC}
\setlength{\parindent}{0pt}
\setlength{\parskip}{0.45em}
\pagestyle{fancy}
\fancyhf{}
\fancyhead[L]{商务智能}
\fancyfoot[C]{\thepage}
\hypersetup{colorlinks=true,linkcolor=black,urlcolor=blue}
\begin{document}
\title{12、间接访问的基本思想}
\author{}
\date{}
\maketitle

数据仓库的间接访问不是让操作型系统实时查询庞大的数据仓库，而是：\par

后台程序周期性分析数据仓库；\par
生成一个较小、易访问的在线预分析数据文件；\par
操作型环境访问这个小文件，而不是直接访问整个数据仓库。\par

两个例子：\par

航空公司佣金计算系统：\par
最佳佣金计算需要当前预定情况和历史订票情况，航空公司离线完成佣金计算和航班历史分析，并创建小型航班状态表，供在线交互时快速查询。\par

零售个性化系统：\par
销售代表与顾客电话咨询时，需要立即查询顾客个性化信息，如上次购物日期、上次购物类型、市场类别信息等。这些信息由数据仓库后台分析程序定时提供。\par

在线预分析文件特点\par

每个数据单元只包含少量数据\par
总体可包含大量信息\par
准确包含在线处理人员需要的内容\par
不被修改\par
批量方式周期性更新\par
是在线高性能访问环境的一部分\par
访问效率高\par
适合访问单个数据单元，而不是大量数据\par

\end{document}
